% preamble.tex — packages and macros
% Include this once via % preamble.tex — packages and macros
% Include this once via % preamble.tex — packages and macros
% Include this once via % preamble.tex — packages and macros
% Include this once via \input{preamble} in main.tex

% ── Encoding & fonts ──────────────────────────────────────────────────────────
\usepackage[utf8]{inputenc}
\usepackage[T1]{fontenc}

% ── Mathematics ───────────────────────────────────────────────────────────────
\usepackage{amsmath, amssymb, amsthm, amsfonts}

% ── Algorithms ────────────────────────────────────────────────────────────────
\usepackage{algorithm}
\usepackage{algpseudocode}

% ── Tables & figures ──────────────────────────────────────────────────────────
\usepackage{booktabs}       % \toprule \midrule \bottomrule
\usepackage{array}          % \raggedright\arraybackslash in p{} columns
\usepackage{graphicx}       % \includegraphics
\usepackage{subcaption}     % subfigures side-by-side

% ── Layout ────────────────────────────────────────────────────────────────────
\usepackage[a4paper, top=0.4in, bottom=0.7in, left=1in, right=1in]{geometry}
\usepackage{microtype}      % better line breaking

% ── Line numbers (blue, padded 001/002/003, dual-side left+right) ─────────────
\usepackage{fmtcount}
\usepackage{lineno}
\usepackage{xcolor}
\definecolor{lineblue}{RGB}{0, 102, 204}
\renewcommand{\linenumberfont}{\normalfont\fontsize{10}{12}\selectfont\color{lineblue}}
\renewcommand{\thelinenumber}{\padzeroes[3]{\decimal{linenumber}}}
\setlength{\linenumbersep}{12pt}
\makeatletter
\def\makeLineNumberLeft{\hss\linenumberfont\thelinenumber\hskip\linenumbersep}
\def\makeLineNumberRight{\linenumberfont\thelinenumber\hss}
\renewcommand\makeLineNumber{%
  \llap{\makeLineNumberLeft}%
  % \rlap{\hskip\textwidth\hskip\linenumbersep\makeLineNumberRight}%
}
\makeatother

% ── Hyperlinks & references ───────────────────────────────────────────────────
\usepackage[colorlinks=true, linkcolor=blue, citecolor=blue, urlcolor=blue]{hyperref}
\usepackage{cleveref}       % \cref{} adds "Theorem", "Figure" etc. automatically

% ── Theorem environments ──────────────────────────────────────────────────────
\newtheorem{theorem}{Theorem}
\newtheorem{lemma}[theorem]{Lemma}
\newtheorem{corollary}[theorem]{Corollary}
\newtheorem{definition}{Definition}
\newtheorem{remark}{Remark}

% ── Canonical notation macros ─────────────────────────────────────────────────
% These match CLAUDE.md exactly. Use these commands everywhere so that
% renaming a symbol only requires changing one line here.

\newcommand{\kw}{k_w}                          % nodes inside window w
\newcommand{\nw}{n_w}                          % testing volume inside window w
\newcommand{\Ktot}{K}                          % total nodes nationally
\newcommand{\Ntot}{N}                          % total testing volume nationally
\newcommand{\cobs}{c}                          % observed cases inside window
\newcommand{\Ctot}{C}                          % total cases nationally
\newcommand{\muexp}{\mu}                       % expected cases = C*n_w/N
\newcommand{\LLR}{\mathrm{LLR}}               % log-likelihood ratio
\newcommand{\RR}{\mathrm{RR}}                 % relative risk
\newcommand{\sens}{\Delta(\kw)}               % sensitivity bound (Theorem 1)
\newcommand{\util}{u(w)}                       % normalized utility = LLR/Delta
\newcommand{\gapu}{\mathrm{gap}_u}            % utility gap between top 2 windows
\newcommand{\Wcal}{\mathcal{W}}               % set of candidate windows
\newcommand{\wstar}{w^*}                       % selected window
\newcommand{\eps}{\varepsilon}                 % privacy budget epsilon
\newcommand{\rhow}{\rho_{\mathrm{week}}}      % per-week zCDP cost
\newcommand{\rhotot}{\rho_{\mathrm{total}}}   % total zCDP cost
\newcommand{\Tmax}{T_{\max}}                  % max surveillance weeks

% ── Draft helpers ─────────────────────────────────────────────────────────────
\newcommand{\todo}[1]{\textcolor{red}{\textbf{[TODO: #1]}}}
\newcommand{\placeholder}[1]{\textcolor{gray}{\textit{[#1]}}}
\newcommand{\needscite}{\textcolor{orange}{[CITE]}}
 in main.tex

% ── Encoding & fonts ──────────────────────────────────────────────────────────
\usepackage[utf8]{inputenc}
\usepackage[T1]{fontenc}

% ── Mathematics ───────────────────────────────────────────────────────────────
\usepackage{amsmath, amssymb, amsthm, amsfonts}

% ── Algorithms ────────────────────────────────────────────────────────────────
\usepackage{algorithm}
\usepackage{algpseudocode}

% ── Tables & figures ──────────────────────────────────────────────────────────
\usepackage{booktabs}       % \toprule \midrule \bottomrule
\usepackage{array}          % \raggedright\arraybackslash in p{} columns
\usepackage{graphicx}       % \includegraphics
\usepackage{subcaption}     % subfigures side-by-side

% ── Layout ────────────────────────────────────────────────────────────────────
\usepackage[a4paper, top=0.4in, bottom=0.7in, left=1in, right=1in]{geometry}
\usepackage{microtype}      % better line breaking

% ── Line numbers (blue, padded 001/002/003, dual-side left+right) ─────────────
\usepackage{fmtcount}
\usepackage{lineno}
\usepackage{xcolor}
\definecolor{lineblue}{RGB}{0, 102, 204}
\renewcommand{\linenumberfont}{\normalfont\fontsize{10}{12}\selectfont\color{lineblue}}
\renewcommand{\thelinenumber}{\padzeroes[3]{\decimal{linenumber}}}
\setlength{\linenumbersep}{12pt}
\makeatletter
\def\makeLineNumberLeft{\hss\linenumberfont\thelinenumber\hskip\linenumbersep}
\def\makeLineNumberRight{\linenumberfont\thelinenumber\hss}
\renewcommand\makeLineNumber{%
  \llap{\makeLineNumberLeft}%
  % \rlap{\hskip\textwidth\hskip\linenumbersep\makeLineNumberRight}%
}
\makeatother

% ── Hyperlinks & references ───────────────────────────────────────────────────
\usepackage[colorlinks=true, linkcolor=blue, citecolor=blue, urlcolor=blue]{hyperref}
\usepackage{cleveref}       % \cref{} adds "Theorem", "Figure" etc. automatically

% ── Theorem environments ──────────────────────────────────────────────────────
\newtheorem{theorem}{Theorem}
\newtheorem{lemma}[theorem]{Lemma}
\newtheorem{corollary}[theorem]{Corollary}
\newtheorem{definition}{Definition}
\newtheorem{remark}{Remark}

% ── Canonical notation macros ─────────────────────────────────────────────────
% These match CLAUDE.md exactly. Use these commands everywhere so that
% renaming a symbol only requires changing one line here.

\newcommand{\kw}{k_w}                          % nodes inside window w
\newcommand{\nw}{n_w}                          % testing volume inside window w
\newcommand{\Ktot}{K}                          % total nodes nationally
\newcommand{\Ntot}{N}                          % total testing volume nationally
\newcommand{\cobs}{c}                          % observed cases inside window
\newcommand{\Ctot}{C}                          % total cases nationally
\newcommand{\muexp}{\mu}                       % expected cases = C*n_w/N
\newcommand{\LLR}{\mathrm{LLR}}               % log-likelihood ratio
\newcommand{\RR}{\mathrm{RR}}                 % relative risk
\newcommand{\sens}{\Delta(\kw)}               % sensitivity bound (Theorem 1)
\newcommand{\util}{u(w)}                       % normalized utility = LLR/Delta
\newcommand{\gapu}{\mathrm{gap}_u}            % utility gap between top 2 windows
\newcommand{\Wcal}{\mathcal{W}}               % set of candidate windows
\newcommand{\wstar}{w^*}                       % selected window
\newcommand{\eps}{\varepsilon}                 % privacy budget epsilon
\newcommand{\rhow}{\rho_{\mathrm{week}}}      % per-week zCDP cost
\newcommand{\rhotot}{\rho_{\mathrm{total}}}   % total zCDP cost
\newcommand{\Tmax}{T_{\max}}                  % max surveillance weeks

% ── Draft helpers ─────────────────────────────────────────────────────────────
\newcommand{\todo}[1]{\textcolor{red}{\textbf{[TODO: #1]}}}
\newcommand{\placeholder}[1]{\textcolor{gray}{\textit{[#1]}}}
\newcommand{\needscite}{\textcolor{orange}{[CITE]}}
 in main.tex

% ── Encoding & fonts ──────────────────────────────────────────────────────────
\usepackage[utf8]{inputenc}
\usepackage[T1]{fontenc}

% ── Mathematics ───────────────────────────────────────────────────────────────
\usepackage{amsmath, amssymb, amsthm, amsfonts}

% ── Algorithms ────────────────────────────────────────────────────────────────
\usepackage{algorithm}
\usepackage{algpseudocode}

% ── Tables & figures ──────────────────────────────────────────────────────────
\usepackage{booktabs}       % \toprule \midrule \bottomrule
\usepackage{array}          % \raggedright\arraybackslash in p{} columns
\usepackage{graphicx}       % \includegraphics
\usepackage{subcaption}     % subfigures side-by-side

% ── Layout ────────────────────────────────────────────────────────────────────
\usepackage[a4paper, top=0.4in, bottom=0.7in, left=1in, right=1in]{geometry}
\usepackage{microtype}      % better line breaking

% ── Line numbers (blue, padded 001/002/003, dual-side left+right) ─────────────
\usepackage{fmtcount}
\usepackage{lineno}
\usepackage{xcolor}
\definecolor{lineblue}{RGB}{0, 102, 204}
\renewcommand{\linenumberfont}{\normalfont\fontsize{10}{12}\selectfont\color{lineblue}}
\renewcommand{\thelinenumber}{\padzeroes[3]{\decimal{linenumber}}}
\setlength{\linenumbersep}{12pt}
\makeatletter
\def\makeLineNumberLeft{\hss\linenumberfont\thelinenumber\hskip\linenumbersep}
\def\makeLineNumberRight{\linenumberfont\thelinenumber\hss}
\renewcommand\makeLineNumber{%
  \llap{\makeLineNumberLeft}%
  % \rlap{\hskip\textwidth\hskip\linenumbersep\makeLineNumberRight}%
}
\makeatother

% ── Hyperlinks & references ───────────────────────────────────────────────────
\usepackage[colorlinks=true, linkcolor=blue, citecolor=blue, urlcolor=blue]{hyperref}
\usepackage{cleveref}       % \cref{} adds "Theorem", "Figure" etc. automatically

% ── Theorem environments ──────────────────────────────────────────────────────
\newtheorem{theorem}{Theorem}
\newtheorem{lemma}[theorem]{Lemma}
\newtheorem{corollary}[theorem]{Corollary}
\newtheorem{definition}{Definition}
\newtheorem{remark}{Remark}

% ── Canonical notation macros ─────────────────────────────────────────────────
% These match CLAUDE.md exactly. Use these commands everywhere so that
% renaming a symbol only requires changing one line here.

\newcommand{\kw}{k_w}                          % nodes inside window w
\newcommand{\nw}{n_w}                          % testing volume inside window w
\newcommand{\Ktot}{K}                          % total nodes nationally
\newcommand{\Ntot}{N}                          % total testing volume nationally
\newcommand{\cobs}{c}                          % observed cases inside window
\newcommand{\Ctot}{C}                          % total cases nationally
\newcommand{\muexp}{\mu}                       % expected cases = C*n_w/N
\newcommand{\LLR}{\mathrm{LLR}}               % log-likelihood ratio
\newcommand{\RR}{\mathrm{RR}}                 % relative risk
\newcommand{\sens}{\Delta(\kw)}               % sensitivity bound (Theorem 1)
\newcommand{\util}{u(w)}                       % normalized utility = LLR/Delta
\newcommand{\gapu}{\mathrm{gap}_u}            % utility gap between top 2 windows
\newcommand{\Wcal}{\mathcal{W}}               % set of candidate windows
\newcommand{\wstar}{w^*}                       % selected window
\newcommand{\eps}{\varepsilon}                 % privacy budget epsilon
\newcommand{\rhow}{\rho_{\mathrm{week}}}      % per-week zCDP cost
\newcommand{\rhotot}{\rho_{\mathrm{total}}}   % total zCDP cost
\newcommand{\Tmax}{T_{\max}}                  % max surveillance weeks

% ── Draft helpers ─────────────────────────────────────────────────────────────
\newcommand{\todo}[1]{\textcolor{red}{\textbf{[TODO: #1]}}}
\newcommand{\placeholder}[1]{\textcolor{gray}{\textit{[#1]}}}
\newcommand{\needscite}{\textcolor{orange}{[CITE]}}
 in main.tex

% ── Encoding & fonts ──────────────────────────────────────────────────────────
\usepackage[utf8]{inputenc}
\usepackage[T1]{fontenc}

% ── Mathematics ───────────────────────────────────────────────────────────────
\usepackage{amsmath, amssymb, amsthm, amsfonts}

% ── Algorithms ────────────────────────────────────────────────────────────────
\usepackage{algorithm}
\usepackage{algpseudocode}

% ── Tables & figures ──────────────────────────────────────────────────────────
\usepackage{booktabs}       % \toprule \midrule \bottomrule
\usepackage{array}          % \raggedright\arraybackslash in p{} columns
\usepackage{graphicx}       % \includegraphics
\usepackage{subcaption}     % subfigures side-by-side

% ── Layout ────────────────────────────────────────────────────────────────────
\usepackage[a4paper, top=0.4in, bottom=0.7in, left=1in, right=1in]{geometry}
\usepackage{microtype}      % better line breaking

% ── Line numbers (blue, padded 001/002/003, dual-side left+right) ─────────────
\usepackage{fmtcount}
\usepackage{lineno}
\usepackage{xcolor}
\definecolor{lineblue}{RGB}{0, 102, 204}
\renewcommand{\linenumberfont}{\normalfont\fontsize{10}{12}\selectfont\color{lineblue}}
\renewcommand{\thelinenumber}{\padzeroes[3]{\decimal{linenumber}}}
\setlength{\linenumbersep}{12pt}
\makeatletter
\def\makeLineNumberLeft{\hss\linenumberfont\thelinenumber\hskip\linenumbersep}
\def\makeLineNumberRight{\linenumberfont\thelinenumber\hss}
\renewcommand\makeLineNumber{%
  \llap{\makeLineNumberLeft}%
  % \rlap{\hskip\textwidth\hskip\linenumbersep\makeLineNumberRight}%
}
\makeatother

% ── Hyperlinks & references ───────────────────────────────────────────────────
\usepackage[colorlinks=true, linkcolor=blue, citecolor=blue, urlcolor=blue]{hyperref}
\usepackage{cleveref}       % \cref{} adds "Theorem", "Figure" etc. automatically

% ── Theorem environments ──────────────────────────────────────────────────────
\newtheorem{theorem}{Theorem}
\newtheorem{lemma}[theorem]{Lemma}
\newtheorem{corollary}[theorem]{Corollary}
\newtheorem{definition}{Definition}
\newtheorem{remark}{Remark}

% ── Canonical notation macros ─────────────────────────────────────────────────
% These match CLAUDE.md exactly. Use these commands everywhere so that
% renaming a symbol only requires changing one line here.

\newcommand{\kw}{k_w}                          % nodes inside window w
\newcommand{\nw}{n_w}                          % testing volume inside window w
\newcommand{\Ktot}{K}                          % total nodes nationally
\newcommand{\Ntot}{N}                          % total testing volume nationally
\newcommand{\cobs}{c}                          % observed cases inside window
\newcommand{\Ctot}{C}                          % total cases nationally
\newcommand{\muexp}{\mu}                       % expected cases = C*n_w/N
\newcommand{\LLR}{\mathrm{LLR}}               % log-likelihood ratio
\newcommand{\RR}{\mathrm{RR}}                 % relative risk
\newcommand{\sens}{\Delta(\kw)}               % sensitivity bound (Theorem 1)
\newcommand{\util}{u(w)}                       % normalized utility = LLR/Delta
\newcommand{\gapu}{\mathrm{gap}_u}            % utility gap between top 2 windows
\newcommand{\Wcal}{\mathcal{W}}               % set of candidate windows
\newcommand{\wstar}{w^*}                       % selected window
\newcommand{\eps}{\varepsilon}                 % privacy budget epsilon
\newcommand{\rhow}{\rho_{\mathrm{week}}}      % per-week zCDP cost
\newcommand{\rhotot}{\rho_{\mathrm{total}}}   % total zCDP cost
\newcommand{\Tmax}{T_{\max}}                  % max surveillance weeks

% ── Draft helpers ─────────────────────────────────────────────────────────────
\newcommand{\todo}[1]{\textcolor{red}{\textbf{[TODO: #1]}}}
\newcommand{\placeholder}[1]{\textcolor{gray}{\textit{[#1]}}}
\newcommand{\needscite}{\textcolor{orange}{[CITE]}}
